\subsection*{一步一步建立直观}

\paragraph{1. 把函数值还原成圆上位置}

在单位圆上，横坐标是余弦，纵坐标是正弦。先把角度调到四分之一圈：此时点在圆顶，纵坐标变化暂时最慢，而横坐标正向左快速变化。三角函数的导数因此可以先理解为圆周速度在两个坐标轴上的投影。



\paragraph{2. 速度为什么总与半径垂直}

匀速转动时，速度朝切线方向，是位置向量逆时针转过直角后的方向。因此位置 $(\cos t,\sin t)$ 对应速度 $(-\sin t,\cos t)$。这里角度采用弧度，单位半径让角度增量和弧长一致；若采用度数，必须补上单位换算因子。



\paragraph{3. 改变半径，检查隐藏的假设}

第三种情形加入径向增长，轨迹成为螺线。此时总速度同时包含切向分量和径向分量，不能再直接用“旋转九十度”表示全部速度。这个对照说明公式所依赖的是固定半径与单位角速度，而不是任意看起来像圆的运动。
